\subsection{A correction to the published tolerance assertion}
\label{sec:gy1972-tolerance-correction}

The question at the end of Vasserstein's talk, occurrence 1972-280.20 and
problem GP-1972-101, asks whether pairwise adjacency implies that a finite
set is a simplex. This is a precise question, but it is not an undisseminated
C or D item. The same assertion was published in
\citet[\S11, Remark~(c), p.~111]{vaserstein1976foundations}, immediately after
the definition of the simplicial scheme. The short announcement
\citep{vaserstein1973foundations} was received on 14 November 1972, one day
after the talk, and gives the same left-translation construction. The
detailed paper was received on 26 February 1975 and published in 1976.

The result below is included as corrective development of an A item.
The label A records publication provenance, not mathematical correctness.
This is a deliberate methodological extension of the usual C/D-only
development filter: a fully checkable correction of an exact published
assertion should not be suppressed or reclassified D merely to fit that
filter. The marks \ours\ and \oursComputed\ identify work performed for
this companion; they do not assert that the counterexample or criterion
has never been noticed before.

\paragraph{The exact model.}
Let \(R\) be an associative ring with identity. Write
\(\mathrm{GL}(R)=\bigcup_{n\geq1}\mathrm{GL}_n(R)\), with stabilization
\(g\mapsto\operatorname{diag}(g,1)\). For a total order \(\prec\) of the
coordinates obtained by a permutation of \(\mathbb N\), let
\(U_\prec\) be the subgroup of matrices differing from the identity in
finitely many entries and satisfying
\[
 u_{ii}=1,\qquad u_{ij}=0\quad\text{if }i\ne j\text{ and }j\prec i.
\]
A finite nonempty set \(X\subset\mathrm{GL}(R)\) is a simplex of the
left-coset model when
\[
                 X\subset gU_\prec
       \quad\text{for some }g\in\mathrm{GL}(R)\text{ and one order }\prec.
\]
This is the definition on P000468 after replacing the notebook's circle
coordinate \(a\) by \(I+a\). Reversing upper and lower triangular
conventions gives the same collection of simplices after coordinate
permutations. The 1976 paper uses right cosets and explains in
\S11, Remark~(b), that inversion identifies its model with the
left-translation model. In particular, the same order must work for an
entire simplex. The pairwise condition permits different orders for
different pairs and is strictly weaker.

\begin{proposition}[Exact acyclic-support criterion]\ours
\label{prop:gy1972-support-criterion}
Let \(X\subset\mathrm{GL}(R)\) be finite and nonempty, and choose
\(x_0\in X\). Form a finite directed graph on a set of coordinates
containing all nonidentity entries of all \(x_0^{-1}x\), \(x\in X\).
Include the edge \(i\to j\), including a loop if \(i=j\), whenever
\[
                (x_0^{-1}x-I)_{ij}\ne0
                \quad\text{for some }x\in X.
\]
Then \(X\) is a simplex of the left-coset model if and only if this graph
is acyclic. The criterion is unchanged by adjoining identity coordinates.
For the right-coset model, replace \(x_0^{-1}x\) by \(xx_0^{-1}\).
\end{proposition}

\begin{proof}
If \(X\subset gU_\prec\), write \(x_0=gu_0\) and \(x=gu_x\).
Then \(x_0^{-1}x=u_0^{-1}u_x\in U_\prec\). Thus every edge points
strictly forward in the same order, which excludes directed cycles and
loops. Conversely, a finite acyclic directed graph has a topological
ordering. Order the finitely many relevant coordinates accordingly and
extend this to a permutation order of \(\mathbb N\). Absence of loops
says that every \(x_0^{-1}x\) has diagonal entries one, and all other
nonzero entries lie strictly above the diagonal in the chosen order.
Consequently every \(x_0^{-1}x\) belongs to \(U_\prec\), giving
\(X\subset x_0U_\prec\). Adding identity coordinates adds isolated
vertices, so cannot create or remove a directed cycle. The right-coset
proof uses \(xx_0^{-1}\) in exactly the same way.
\end{proof}

This is an exact replacement for the proposed pairwise test. It reduces
simplex recognition to finding a common order of the union of the
relative supports, rather than checking each pair with a separate order.
When equality to zero is decidable in the coefficient ring, the criterion
also gives a finite decision procedure: compute the relative matrices
and perform a topological-sort test on their support graph.

\begin{proposition}[A three-vertex counterexample over every nonzero unital ring]\ours
\label{prop:gy1972-nonflag-triple}
Suppose \(1_R\ne0\), and in \(\mathrm{GL}_3(R)\) put
\[
 a=I+E_{12},\qquad b=I+E_{23},\qquad c=I+E_{31}.
\]
Every two-element subset of \(\{a,b,c\}\) is a simplex, but the entire
set is not a simplex. This holds in both the left-coset and right-coset
conventions and after every stabilization.
\end{proposition}

\begin{proof}
The matrix units carry coefficient \(1_R\), so each is square-zero and
the three inverses are obtained by replacing plus with minus. Direct
multiplication gives
\[
\begin{aligned}
 a^{-1}b&=I-E_{12}+E_{23}-E_{13},\\
 a^{-1}c&=I-E_{12}+E_{31},\\
 b^{-1}c&=I-E_{23}+E_{31}-E_{21}.
\end{aligned}
\]
The off-diagonal supports point forward, respectively, in the orders
\[
             1\prec2\prec3,\qquad
             3\prec1\prec2,\qquad
             2\prec3\prec1.
\]
Thus every pair is a left-coset simplex. If the triple were a simplex,
the same order would triangularize both \(a^{-1}b\) and \(a^{-1}c\).
Their supports would require simultaneously
\(1\prec2\), \(2\prec3\), and \(3\prec1\), an impossibility.
Equivalently, their union support contains a directed triangle, so
Proposition~\ref{prop:gy1972-support-criterion} applies.

For completeness, the same three vertices also work directly in the
right-coset convention, without replacing them by their inverses:
\[
\begin{aligned}
 ba^{-1}&=I-E_{12}+E_{23},\\
 ca^{-1}&=I-E_{12}+E_{31}-E_{32},\\
 cb^{-1}&=I-E_{23}+E_{31}.
\end{aligned}
\]
The same respective three orders triangularize these pairwise
quotients. The first two again have union support containing
\(1\to2\to3\to1\), excluding a common right coset.

All decisive coefficients are \(1_R\) or \(-1_R\), which are nonzero
in every nonzero unital ring, including characteristic two. The proof
uses no field hypothesis and no division by a coefficient. Additional
identity coordinates cannot remove a cycle among coordinates one,
two, and three. This proves all the stated cases.
\end{proof}

\begin{proposition}[Minimal coordinate rank of the obstruction]\ours
\label{prop:gy1972-rank-minimal}
For every unital ring, a finite pairwise adjacent set of vertices
supported in one two-by-two block is a simplex, even when adjacency is
tested after stabilization. Consequently the example in
Proposition~\ref{prop:gy1972-nonflag-triple} has minimal coordinate rank,
as well as the smallest possible number of vertices for failure of the
pairwise criterion.
\end{proposition}

\begin{proof}
Translate by the inverse of one vertex so that the set contains \(I\).
In the left model, every other vertex adjacent to \(I\) and supported
on two coordinates must have one of the forms
\(I+rE_{12}\) or \(I+sE_{21}\). Stabilization does not change this
fact, since only the relative order of those two coordinates matters.
If both orientations occur nontrivially, with \(r\ne0\) and \(s\ne0\),
then
\[
 (I+rE_{12})^{-1}(I+sE_{21})
       =I-rE_{12}+sE_{21}-rsE_{11}.
\]
Both off-diagonal entries are nonzero and point in opposite directions.
No coordinate order can make this matrix unitriangular, irrespective of
whether \(rs\) vanishes. This contradicts pairwise adjacency.
Hence all nonidentity vertices use one orientation and lie in one
unitriangular group, so the whole set is a simplex. The same argument
works with right quotients, or follows by inversion. Rank one is
immediate. Sets of at most two vertices cannot contradict a condition
defined by their two-element subsets, while the displayed triple does.
\end{proof}

\paragraph{Independent verification.}\oursComputed
The included Python program \texttt{check\_tolerance.py} uses exact
matrix arithmetic and exhausts all coordinate permutations in dimensions
three, four, and five over \(\mathbb Z\) and
\(\mathbb F_2,\mathbb F_3,\mathbb F_5,\mathbb F_7\), in both coset
conventions. Its thirty parameter cases all have three present edges
and no triple simplex. The complete successful-order lists, integer
quotient matrices, and results are in \texttt{tolerance\_output.json}.
The integer calculation rechecks the signs and cross terms; the
characteristic-two calculation checks that the obstruction does not
rely on distinguishing plus from minus. A separate independently
written audit by the assembling researcher exhausts the four finite
fields in the same three dimensions and both conventions and agrees in
all twenty-four cases. The general-ring and arbitrary-stabilization
conclusions follow from the proof, not extrapolation from these finite
computations. The rank-two argument was also checked independently by
normalizing a base vertex and recomputing the opposite-direction
quotient above.

\paragraph{Scope and historical assessment.}
This resolves the literal notebook question negatively; no residual
mathematical obstruction remains for that question. The 1976 source
expressly says the tolerance assertion is not used in the paper.
Accordingly, this local correction provides no basis for declaring its
comparison theorems or the Volodin construction invalid. Searching the
speaker's Foundations papers and bibliography, the original Volodin
paper, and targeted literature queries identified the exact published
assertion but did not establish the history of any later correction.
Absence of a retrieved erratum is not evidence that none exists, and no
claim is made about why the error persisted.

The strongest general statement supported after correction is the
acyclic-support criterion above, with its full common-order requirement.
Replacing the simplicial complex by its flag completion would create a
different construction and does not rescue the literal assertion.
A further bounded bibliographic run could trace corrections or uses of
Remark~(c); it is not needed to establish the explicit counterexample,
and it is not recorded as a manufactured open mathematical problem.
